\documentclass[aspectratio=169,10pt]{beamer}
\usepackage{fontspec}
\usepackage{graphicx}
\usepackage{tikz}
\usepackage{amsmath}
\usepackage{unicode-math}
\usetikzlibrary{arrows.meta,backgrounds,calc,fit,positioning,shapes.geometric}

\setmainfont{TeX Gyre Pagella}
\setsansfont{TeX Gyre Heros}
\setmathfont{TeX Gyre Pagella Math}

\definecolor{BMEIvory}{HTML}{F7F2E8}
\definecolor{BMEInk}{HTML}{2F3437}
\definecolor{BMEBrown}{HTML}{5C433B}
\definecolor{BMETerracotta}{HTML}{B66A50}
\definecolor{BMESage}{HTML}{7D9276}
\definecolor{BMESand}{HTML}{E7DAC7}
\definecolor{BMEStone}{HTML}{C9BBA8}
\definecolor{BMEMuted}{HTML}{77716B}
\definecolor{BMEWhite}{HTML}{FFFDFC}
\definecolor{BMEFlow}{HTML}{655A54}
\definecolor{BMEReturn}{HTML}{718C80}
\definecolor{BMEOptional}{HTML}{938B83}

\setbeamercolor{normal text}{fg=BMEInk,bg=BMEIvory}
\setbeamercolor{structure}{fg=BMEBrown}
\setbeamertemplate{navigation symbols}{}
\setbeamertemplate{footline}{}
\setbeamertemplate{headline}{}
\AtBeginEnvironment{frame}{\raggedright\hyphenpenalty=10000\exhyphenpenalty=10000\relax}

\newcommand{\WarmFrameHeader}[3]{%
  \node[anchor=north west,inner sep=0pt,text=BMETerracotta,
        font=\sffamily\fontsize{6.2}{7.5}\selectfont\bfseries]
    at ([xshift=0.86cm,yshift=-0.30cm]current page.north west) {\MakeUppercase{#2}};
  \node[anchor=north west,inner sep=0pt,text=BMEBrown,
        font=\rmfamily\fontsize{16.2}{19.2}\selectfont\bfseries]
    at ([xshift=0.86cm,yshift=-0.58cm]current page.north west)
    {\makebox[11.9cm][l]{#1}};
  \draw[BMETerracotta,line width=1.25pt]
    ([xshift=0.86cm,yshift=-1.29cm]current page.north west) -- ++(0.74cm,0);
  \draw[BMEStone,line width=0.45pt]
    ([xshift=1.72cm,yshift=-1.29cm]current page.north west) --
    ([xshift=-0.86cm,yshift=-1.29cm]current page.north east);
  \node[anchor=north east,inner sep=0pt,text=BMEBrown,
        font=\sffamily\fontsize{8.5}{10}\selectfont\bfseries]
    at ([xshift=-0.86cm,yshift=-0.58cm]current page.north east) {#3};
}

\newcommand{\WarmFrameFooter}[1]{%
  \node[anchor=south west,inner sep=0pt,text=BMEMuted,
        font=\sffamily\fontsize{5.8}{7}\selectfont]
    at ([xshift=0.86cm,yshift=0.20cm]current page.south west) {#1};
  \fill[BMETerracotta]
    ([xshift=-1.02cm,yshift=0.24cm]current page.south east) circle[radius=1.2pt];
}

\newcommand{\WarmTag}[2][BMESand]{%
  \tikz[baseline=(warmtag.base)]{%
    \node[rounded corners=3pt,fill=#1,draw=none,minimum width=0pt,minimum height=0pt,
          outer sep=0pt,inner xsep=3.6pt,inner ysep=1.2pt,text height=1.15ex,text depth=0.20ex,
          text=BMEBrown,font=\sffamily\fontsize{4.9}{5.8}\selectfont\bfseries] (warmtag) {#2};}%
}

\newcommand{\WarmBodyText}{%
  \color{BMEInk}\sffamily\mdseries\fontsize{6.8}{8.8}\selectfont
  \raggedright\rightskip=0pt plus 1fil\spaceskip=0.30em plus 0.05em minus 0.03em\xspaceskip=0.45em\relax}
\newcommand{\WarmBodySmall}{%
  \color{BMEInk}\sffamily\mdseries\fontsize{6.5}{8.3}\selectfont
  \raggedright\rightskip=0pt plus 1fil\spaceskip=0.30em plus 0.05em minus 0.03em\xspaceskip=0.45em\relax}
\newcommand{\WarmDisplayMath}{%
  \color{BMEBrown}\rmfamily\fontsize{13.5}{16.2}\selectfont}

\tikzset{
  slide/prose/.style={align=left,execute at begin node={\raggedright\hyphenpenalty=10000\exhyphenpenalty=10000\relax}},
  warm/panel/.style={rounded corners=3.5pt,draw=BMEStone,line width=0.45pt,fill=BMEWhite},
  warm/node/.style={rounded corners=2.8pt,draw=BMEStone,line width=0.55pt,fill=BMEWhite,
    minimum height=0.72cm,inner xsep=0.18cm,inner ysep=0.12cm,align=center,
    text=BMEInk,font=\sffamily\fontsize{6.5}{7.8}\selectfont},
  warm/action/.style={warm/node,align=left,inner xsep=0.22cm,inner ysep=0.16cm},
  warm/decision/.style={diamond,aspect=1.55,draw=BMEBrown,line width=0.62pt,fill=BMEWhite,
    minimum width=1.18cm,minimum height=0.76cm,inner sep=0.05cm,align=center,
    text=BMEBrown,font=\sffamily\fontsize{5.8}{6.9}\selectfont\bfseries},
  warm/terminal/.style={rounded corners=8pt,draw=BMEStone,line width=0.55pt,fill=BMEWhite,
    minimum height=0.58cm,inner xsep=0.22cm,inner ysep=0.10cm,align=center,
    text=BMEBrown,font=\sffamily\fontsize{6.0}{7.2}\selectfont\bfseries},
  warm/graph edge/.style={draw=BMEStone,line width=1.08pt,line cap=round,line join=round},
  warm/graph node/.style={circle,draw=BMEBrown,line width=0.75pt,fill=BMEWhite,
    minimum size=0.68cm,inner sep=0pt,text=BMEBrown,font=\sffamily\fontsize{9}{10}\selectfont\bfseries},
  warm/graph current/.style={warm/graph node,draw=BMETerracotta,line width=0.95pt,fill=BMESand},
  warm/graph route/.style={draw=BMETerracotta,line width=1.25pt,line cap=round,line join=round,
    shorten <=3.4mm,shorten >=3.4mm},
  warm/graph route arrow/.style={warm/graph route,-{Stealth[length=1.35mm,width=1.0mm]}},
  warm/flow/.style={-{Stealth[length=1.35mm,width=1.0mm]},draw=BMEFlow,line width=0.66pt,
    line cap=round,line join=round},
  warm/return/.style={-{Stealth[length=1.30mm,width=0.95mm]},draw=BMEReturn,line width=0.61pt,
    line cap=round,line join=round},
  warm/optional/.style={-{Stealth[length=1.25mm,width=0.92mm]},draw=BMEOptional,line width=0.56pt,
    dashed,dash pattern=on 3.1pt off 2.5pt,line cap=round,line join=round},
  warm/highlight/.style={-{Stealth[length=1.35mm,width=1.0mm]},draw=BMETerracotta,line width=0.78pt,
    line cap=round,line join=round},
  warm/arrow/.style={warm/flow},
  warm/edge label/.style={text=BMEMuted,font=\sffamily\fontsize{5.4}{6.6}\selectfont,fill=BMEIvory,
    rounded corners=1.5pt,inner xsep=2.4pt,inner ysep=1.1pt},
  warm/adjacent label/.style={text=BMEMuted,font=\sffamily\fontsize{5.4}{6.6}\selectfont,inner sep=0pt},
  warm/label/.style={warm/edge label}
}


\begin{document}

\begin{frame}[plain]
  \begin{tikzpicture}[remember picture,overlay]
    \fill[BMEIvory] (current page.south west) rectangle (current page.north east);
    \fill[BMETerracotta]
      ([xshift=0.92cm,yshift=-1.24cm]current page.north west) rectangle ++(0.14cm,-0.86cm);
    \node[anchor=north west,inner sep=0pt,text=BMETerracotta,
          font=\sffamily\fontsize{6.5}{8}\selectfont\bfseries]
      at ([xshift=1.28cm,yshift=-1.27cm]current page.north west)
      {WORKED EXAMPLE · MATHEMATICAL EXPOSITION};
    \node[anchor=north west,inner sep=0pt,text=BMEBrown,
          font=\rmfamily\fontsize{25}{29}\selectfont\bfseries,text width=13.1cm]
      at ([xshift=1.22cm,yshift=-2.62cm]current page.north west)
      {Random walks on graphs};
    \node[anchor=north west,inner sep=0pt,text=BMEMuted,
          font=\sffamily\fontsize{9.8}{12.5}\selectfont,text width=10.2cm]
      at ([xshift=1.25cm,yshift=-3.80cm]current page.north west)
      {From local movement to graph structure and long-run behavior};
    \draw[BMEStone,line width=0.55pt]
      ([xshift=1.25cm,yshift=-4.72cm]current page.north west) -- ++(6.4cm,0);
    \node[anchor=south west,inner sep=0pt,text=BMEMuted,
          font=\sffamily\fontsize{6.9}{8.6}\selectfont]
      at ([xshift=1.25cm,yshift=0.72cm]current page.south west)
      {Warm Editorial style · seven-slide example};
    \fill[BMESage]
      ([xshift=-1.12cm,yshift=0.82cm]current page.south east) circle[radius=2.1pt];
  \end{tikzpicture}
\end{frame}

\begin{frame}[plain]
  \begin{tikzpicture}[remember picture,overlay]
    \fill[BMEIvory] (current page.south west) rectangle (current page.north east);
    \WarmFrameHeader{Edges become transition probabilities}{FOUNDATION}{02}

    \coordinate (A) at ([xshift=2.05cm,yshift=-3.05cm]current page.north west);
    \coordinate (B) at ([xshift=4.75cm,yshift=-2.18cm]current page.north west);
    \coordinate (C) at ([xshift=7.55cm,yshift=-3.30cm]current page.north west);
    \coordinate (D) at ([xshift=5.32cm,yshift=-5.42cm]current page.north west);
    \coordinate (E) at ([xshift=2.18cm,yshift=-5.18cm]current page.north west);
    \draw[warm/graph edge] (A)--(B)--(C)--(D)--(E)--(A);
    \draw[warm/graph edge] (A)--(D);
    \draw[warm/graph edge] (B)--(D);
    \node[warm/graph node] (nodeA) at (A) {a};
    \node[warm/graph current] (nodeB) at (B) {b};
    \node[warm/graph node] (nodeC) at (C) {c};
    \node[warm/graph node] (nodeD) at (D) {d};
    \node[warm/graph node] (nodeE) at (E) {e};
    \node[anchor=south,inner sep=0pt]
      at ([yshift=0.16cm]nodeB.north) {\WarmTag[BMESand]{CURRENT}};
    \draw[warm/highlight] (nodeB.south east) to[out=-58,in=72] (nodeD.north east);
    \node[warm/adjacent label,anchor=west,text=BMETerracotta,font=\sffamily\fontsize{5.8}{7}\selectfont\bfseries]
      at ([xshift=0.75cm,yshift=-1.25cm]nodeB.east) {possible step};

    \draw[BMEStone,line width=0.45pt]
      ([xshift=8.55cm,yshift=-1.62cm]current page.north west) --
      ([xshift=8.55cm,yshift=-7.55cm]current page.north west);
    \node[anchor=north west,inner sep=0pt,slide/prose,text width=5.78cm]
      at ([xshift=9.05cm,yshift=-1.82cm]current page.north west)
      {\color{BMETerracotta}\sffamily\fontsize{6.1}{7.4}\selectfont\bfseries DEFINITION\par
       \vspace{0.14cm}
       \color{BMEBrown}\rmfamily\fontsize{10.4}{12.4}\selectfont\bfseries
       Choose a neighboring vertex in proportion to edge weight.\par
       \vspace{0.26cm}
       \WarmBodyText
       For weighted adjacency matrix $W$, the transition probability is\par
       \vspace{0.22cm}
       \WarmDisplayMath
       \[
         P_{ij}=\frac{W_{ij}}{\sum_k W_{ik}}.
       \]
       \vspace{0.08cm}
       \WarmBodyText
       Every row of $P$ sums to one. The graph specifies the available transitions.\par};

    \WarmFrameFooter{Random walks · local rule}
  \end{tikzpicture}
\end{frame}

\begin{frame}[plain]
  \begin{tikzpicture}[remember picture,overlay]
    \fill[BMEIvory] (current page.south west) rectangle (current page.north east);
    \WarmFrameHeader{One step is matrix multiplication}{DYNAMICS}{03}

    \node[anchor=north west,inner sep=0pt,slide/prose,text width=6.55cm]
      at ([xshift=0.98cm,yshift=-1.78cm]current page.north west)
      {\color{BMETerracotta}\sffamily\fontsize{6.1}{7.4}\selectfont\bfseries TRANSITION MATRIX\par
       \vspace{0.22cm}
       \color{BMEBrown}\rmfamily\fontsize{11.0}{13.4}\selectfont
       \[
       P=\begin{pmatrix}
       0 & \tfrac13 & 0 & \tfrac13 & \tfrac13\\
       \tfrac13 & 0 & \tfrac13 & \tfrac13 & 0\\
       0 & \tfrac12 & 0 & \tfrac12 & 0\\
       \tfrac14 & \tfrac14 & \tfrac14 & 0 & \tfrac14\\
       \tfrac12 & 0 & 0 & \tfrac12 & 0
       \end{pmatrix}.
       \]
       \vspace{0.18cm}
       \WarmBodyText
       At each vertex, choose one neighbor uniformly.\par};

    \draw[BMEStone,line width=0.45pt]
      ([xshift=8.28cm,yshift=-1.62cm]current page.north west) --
      ([xshift=8.28cm,yshift=-7.55cm]current page.north west);

    \node[anchor=north west,inner sep=0pt,slide/prose,text width=6.18cm]
      at ([xshift=8.80cm,yshift=-1.78cm]current page.north west)
      {\color{BMESage}\sffamily\fontsize{6.1}{7.4}\selectfont\bfseries DISTRIBUTION UPDATE\par
       \vspace{0.20cm}
       \color{BMEBrown}\rmfamily\fontsize{10.2}{12.4}\selectfont\bfseries
       If $\mu_t$ records the position distribution at time $t$, then\par
       \vspace{0.18cm}
       \WarmDisplayMath
       \[
         \mu_{t+1}=\mu_tP.
       \]
       \vspace{0.18cm}
       \WarmBodyText
       The initial row vector is concentrated at vertex $b$.\par
       \vspace{0.18cm}
       After one step, each of $a$, $c$, and $d$ has probability $1/3$.\par};

    \WarmFrameFooter{Random walks · linear evolution}
  \end{tikzpicture}
\end{frame}

\begin{frame}[plain]
  \begin{tikzpicture}[remember picture,overlay]
    \fill[BMEIvory] (current page.south west) rectangle (current page.north east);
    \WarmFrameHeader{A bottleneck delays movement between two regions}{STRUCTURE}{04}

    \fill[BMEWhite,rounded corners=5pt]
      ([xshift=0.88cm,yshift=-1.60cm]current page.north west) rectangle
      ([xshift=5.12cm,yshift=-6.88cm]current page.north west);
    \fill[BMEWhite,rounded corners=5pt]
      ([xshift=5.60cm,yshift=-1.60cm]current page.north west) rectangle
      ([xshift=9.92cm,yshift=-6.88cm]current page.north west);
    \node[anchor=north west,inner sep=0pt,text=BMETerracotta,
          font=\sffamily\fontsize{5.8}{7}\selectfont\bfseries]
      at ([xshift=1.10cm,yshift=-1.82cm]current page.north west) {REGION A};
    \node[anchor=north west,inner sep=0pt,text=BMESage,
          font=\sffamily\fontsize{5.8}{7}\selectfont\bfseries]
      at ([xshift=5.82cm,yshift=-1.82cm]current page.north west) {REGION B};

    \coordinate (GA) at ([xshift=1.70cm,yshift=-3.02cm]current page.north west);
    \coordinate (GB) at ([xshift=3.38cm,yshift=-2.58cm]current page.north west);
    \coordinate (GC) at ([xshift=2.12cm,yshift=-4.70cm]current page.north west);
    \coordinate (GD) at ([xshift=4.28cm,yshift=-4.10cm]current page.north west);
    \coordinate (GE) at ([xshift=6.48cm,yshift=-4.10cm]current page.north west);
    \coordinate (GF) at ([xshift=7.22cm,yshift=-2.62cm]current page.north west);
    \coordinate (GG) at ([xshift=9.08cm,yshift=-3.12cm]current page.north west);
    \coordinate (GH) at ([xshift=8.34cm,yshift=-5.02cm]current page.north west);

    \draw[warm/graph edge] (GA)--(GB)--(GD)--(GC)--(GA);
    \draw[warm/graph edge] (GB)--(GC);
    \draw[warm/graph edge] (GE)--(GF)--(GG)--(GH)--(GE);
    \draw[warm/graph edge] (GE)--(GG);
    \draw[warm/graph edge] (GF)--(GH);
    \draw[BMEOptional,line width=0.82pt,dashed,dash pattern=on 3pt off 2.3pt,
          line cap=round] (GD)--(GE);

    \node[warm/graph current] (graph-a) at (GA) {a};
    \node[warm/graph node] (graph-b) at (GB) {b};
    \node[warm/graph node] (graph-c) at (GC) {c};
    \node[warm/graph node] (graph-d) at (GD) {d};
    \node[warm/graph node] (graph-e) at (GE) {e};
    \node[warm/graph node] (graph-f) at (GF) {f};
    \node[warm/graph current,draw=BMESage,fill=BMEIvory] (graph-g) at (GG) {g};
    \node[warm/graph node] (graph-h) at (GH) {h};

    \draw[warm/graph route] (graph-a.center)--(graph-c.center);
    \draw[warm/graph route] (graph-c.center)--(graph-d.center);
    \draw[warm/graph route arrow] (graph-d.center)--(graph-e.center);
    \draw[warm/graph route] (graph-e.center)--(graph-h.center);
    \draw[warm/graph route] (graph-h.center)--(graph-g.center);

    \node[anchor=south,inner sep=0pt]
      at ([yshift=0.15cm]graph-a.north) {\WarmTag[BMESand]{CURRENT}};
    \node[anchor=south,inner sep=0pt]
      at ([yshift=0.15cm]graph-g.north) {\WarmTag[BMEIvory]{TARGET}};
    \node[warm/adjacent label,anchor=south,text=BMEMuted]
      at ([yshift=0.13cm]$(graph-d)!0.5!(graph-e)$) {bridge weight $0.1$};

    \draw[BMEStone,line width=0.45pt]
      ([xshift=10.35cm,yshift=-1.62cm]current page.north west) --
      ([xshift=10.35cm,yshift=-7.55cm]current page.north west);
    \node[anchor=north west,inner sep=0pt,slide/prose,text width=4.22cm]
      at ([xshift=10.80cm,yshift=-1.82cm]current page.north west)
      {\color{BMETerracotta}\sffamily\fontsize{6.1}{7.4}\selectfont\bfseries CUT BOTTLENECK\par
       \vspace{0.15cm}
       \color{BMEBrown}\rmfamily\fontsize{10.2}{12.2}\selectfont\bfseries
       A single weak edge connects the two regions.\par
       \vspace{0.24cm}
       \WarmDisplayMath
       \[
         \Phi(S)=\frac{w(S,\bar S)}{\operatorname{vol}(S)}.
       \]
       \vspace{0.08cm}
       \WarmBodyText
       Small conductance means that probability leaves a region slowly.\par
       \vspace{0.19cm}
       Crossing remains possible but rare.\par};

    \WarmFrameFooter{Random walks · graph structure and conductance}
  \end{tikzpicture}
\end{frame}

\begin{frame}[plain]
  \begin{tikzpicture}[remember picture,overlay]
    \fill[BMEIvory] (current page.south west) rectangle (current page.north east);
    \WarmFrameHeader{Walk variants produce different convergence profiles}{COMPARISON}{05}

    \fill[BMEWhite,rounded corners=5pt]
      ([xshift=0.88cm,yshift=-1.60cm]current page.north west) rectangle
      ([xshift=10.95cm,yshift=-7.18cm]current page.north west);
    \node[anchor=north west,inner sep=0pt,text=BMETerracotta,
          font=\sffamily\fontsize{5.9}{7.1}\selectfont\bfseries]
      at ([xshift=1.14cm,yshift=-1.84cm]current page.north west)
      {TOTAL VARIATION DISTANCE};

    \fill[BMESand,opacity=0.38]
      ([xshift=1.48cm,yshift=-5.60cm]current page.north west) rectangle
      ([xshift=9.80cm,yshift=-6.30cm]current page.north west);
    \draw[BMEStone,line width=0.65pt]
      ([xshift=1.48cm,yshift=-6.30cm]current page.north west) --
      ([xshift=9.82cm,yshift=-6.30cm]current page.north west);
    \draw[BMEStone,line width=0.65pt]
      ([xshift=1.48cm,yshift=-6.30cm]current page.north west) --
      ([xshift=1.48cm,yshift=-2.30cm]current page.north west);
    \foreach \x/\label in {1.48/0,3.12/2,4.76/4,6.40/6,8.04/8,9.68/10}{
      \draw[BMEStone,line width=0.4pt]
        ([xshift=\x cm,yshift=-6.24cm]current page.north west) -- ++(0,-0.11cm);
      \node[anchor=north,inner sep=0pt,text=BMEMuted,font=\sffamily\fontsize{5.2}{6.3}\selectfont]
        at ([xshift=\x cm,yshift=-6.46cm]current page.north west) {\label};
    }
    \foreach \y/\label in {-6.30/0,-5.34/.25,-4.38/.50,-3.42/.75,-2.46/1.0}{
      \draw[BMEStone,line width=0.4pt]
        ([xshift=1.42cm,yshift=\y cm]current page.north west) -- ++(-0.11cm,0);
      \node[anchor=east,inner sep=0pt,text=BMEMuted,font=\sffamily\fontsize{5.2}{6.3}\selectfont]
        at ([xshift=1.24cm,yshift=\y cm]current page.north west) {\label};
    }
    \draw[BMEOptional,line width=0.52pt,dashed]
      ([xshift=1.48cm,yshift=-5.60cm]current page.north west) --
      ([xshift=9.80cm,yshift=-5.60cm]current page.north west);
    \node[anchor=south west,inner sep=0pt,text=BMEMuted,
          font=\sffamily\fontsize{5.2}{6.3}\selectfont]
      at ([xshift=1.64cm,yshift=-5.52cm]current page.north west) {$0.15$ threshold};

    \draw[BMETerracotta,line width=1.25pt,line cap=round]
      ([xshift=1.48cm,yshift=-2.55cm]current page.north west)
      .. controls ([xshift=2.45cm,yshift=-3.86cm]current page.north west) and
                  ([xshift=5.30cm,yshift=-5.20cm]current page.north west) ..
      ([xshift=9.20cm,yshift=-5.78cm]current page.north west);
    \draw[BMESage,line width=1.25pt,line cap=round]
      ([xshift=1.48cm,yshift=-2.55cm]current page.north west)
      .. controls ([xshift=3.10cm,yshift=-3.08cm]current page.north west) and
                  ([xshift=6.25cm,yshift=-4.18cm]current page.north west) ..
      ([xshift=9.20cm,yshift=-4.96cm]current page.north west);
    \draw[BMEBrown,line width=1.05pt,dashed,dash pattern=on 4pt off 2.2pt,line cap=round]
      ([xshift=1.48cm,yshift=-2.55cm]current page.north west)
      .. controls ([xshift=2.22cm,yshift=-4.22cm]current page.north west) and
                  ([xshift=4.75cm,yshift=-5.28cm]current page.north west) ..
      ([xshift=9.20cm,yshift=-5.42cm]current page.north west);
    \node[anchor=west,inner sep=0pt,text=BMETerracotta,font=\sffamily\fontsize{5.6}{6.8}\selectfont\bfseries]
      at ([xshift=9.38cm,yshift=-5.78cm]current page.north west) {simple};
    \node[anchor=west,inner sep=0pt,text=BMEBrown,font=\sffamily\fontsize{5.6}{6.8}\selectfont\bfseries]
      at ([xshift=9.38cm,yshift=-5.42cm]current page.north west) {restart};
    \node[anchor=west,inner sep=0pt,text=BMESage,font=\sffamily\fontsize{5.6}{6.8}\selectfont\bfseries]
      at ([xshift=9.38cm,yshift=-4.96cm]current page.north west) {lazy};
    \node[anchor=north,inner sep=0pt,text=BMEMuted,font=\sffamily\fontsize{5.7}{6.9}\selectfont]
      at ([xshift=5.68cm,yshift=-6.82cm]current page.north west) {time step $t$};

    \node[warm/panel,slide/prose,anchor=north west,inner sep=0.25cm,text width=3.50cm,minimum height=4.84cm]
      at ([xshift=11.34cm,yshift=-1.64cm]current page.north west)
      {\color{BMETerracotta}\sffamily\fontsize{6.0}{7.2}\selectfont\bfseries ILLUSTRATIVE READOUT\par
       \vspace{0.20cm}
       \color{BMEBrown}\rmfamily\fontsize{9.5}{11.4}\selectfont\bfseries
       Fast early decline can end at a biased floor.\par
       \vspace{0.25cm}
       \WarmBodySmall
       \textbf{Simple walk}\par threshold near $t=8$\par
       \vspace{0.18cm}
       \textbf{Restart walk}\par reaches a biased floor\par
       \vspace{0.18cm}
       \textbf{Lazy walk}\par slower, same target\par
       \vspace{0.24cm}
       Values are schematic only.\par};

    \WarmFrameFooter{Schematic profiles · direct labels and threshold band}
  \end{tikzpicture}
\end{frame}

\begin{frame}[plain]
  \begin{tikzpicture}[remember picture,overlay]
    \fill[BMEIvory] (current page.south west) rectangle (current page.north east);
    \WarmFrameHeader{This walk converges to one stationary distribution}{LONG-RUN BEHAVIOR}{06}

    \node[anchor=north west,inner sep=0pt,text=BMETerracotta,
          font=\sffamily\fontsize{6.1}{7.4}\selectfont\bfseries]
      at ([xshift=0.98cm,yshift=-1.66cm]current page.north west) {DISTANCE TO STATIONARITY};
    \draw[BMEStone,line width=0.65pt]
      ([xshift=1.42cm,yshift=-6.32cm]current page.north west) --
      ([xshift=8.20cm,yshift=-6.32cm]current page.north west);
    \draw[BMEStone,line width=0.65pt]
      ([xshift=1.42cm,yshift=-6.32cm]current page.north west) --
      ([xshift=1.42cm,yshift=-2.08cm]current page.north west);
    \foreach \x/\label in {1.42/0,2.77/2,4.12/4,5.47/6,6.82/8,8.17/10}{
      \draw[BMEStone,line width=0.4pt]
        ([xshift=\x cm,yshift=-6.26cm]current page.north west) -- ++(0,-0.11cm);
      \node[anchor=north,inner sep=0pt,text=BMEMuted,font=\sffamily\fontsize{5.4}{6.6}\selectfont]
        at ([xshift=\x cm,yshift=-6.48cm]current page.north west) {\label};
    }
    \foreach \y/\label in {-5.48/.2,-4.63/.4,-3.78/.6,-2.93/.8}{
      \draw[BMEStone,line width=0.4pt]
        ([xshift=1.36cm,yshift=\y cm]current page.north west) -- ++(-0.11cm,0);
      \node[anchor=east,inner sep=0pt,text=BMEMuted,font=\sffamily\fontsize{5.4}{6.6}\selectfont]
        at ([xshift=1.17cm,yshift=\y cm]current page.north west) {\label};
    }
    \draw[BMETerracotta,line width=1.35pt]
      ([xshift=1.42cm,yshift=-2.55cm]current page.north west)
      .. controls ([xshift=2.25cm,yshift=-4.15cm]current page.north west) and
                  ([xshift=4.20cm,yshift=-5.42cm]current page.north west) ..
      ([xshift=8.15cm,yshift=-6.05cm]current page.north west);
    \node[anchor=north,inner sep=0pt,text=BMEMuted,font=\sffamily\fontsize{5.8}{7}\selectfont]
      at ([xshift=4.82cm,yshift=-6.82cm]current page.north west) {time step $t$};

    \draw[BMEStone,line width=0.45pt]
      ([xshift=8.80cm,yshift=-1.62cm]current page.north west) --
      ([xshift=8.80cm,yshift=-7.55cm]current page.north west);
    \node[anchor=north west,inner sep=0pt,slide/prose,text width=5.55cm]
      at ([xshift=9.28cm,yshift=-1.80cm]current page.north west)
      {\color{BMETerracotta}\sffamily\fontsize{6.1}{7.4}\selectfont\bfseries STATIONARY DISTRIBUTION\par
       \vspace{0.18cm}
       \WarmDisplayMath
       \[
         \pi P=\pi.
       \]
       \vspace{0.12cm}
       \WarmBodyText
       On a finite, connected, non-bipartite graph, the walk has a unique limit.\par
       \vspace{0.25cm}
       \color{BMEBrown}\rmfamily\itshape\fontsize{8.8}{10.8}\selectfont
       Vertices of higher degree receive more long-run probability.\par};

    \WarmFrameFooter{Schematic curve · conceptual illustration}
  \end{tikzpicture}
\end{frame}

\begin{frame}[plain]
  \begin{tikzpicture}[remember picture,overlay]
    \fill[BMEIvory] (current page.south west) rectangle (current page.north east);
    \WarmFrameHeader{Applications, assumptions, and reporting choices}{INTERPRETATION}{07}

    \node[anchor=north west,inner sep=0pt,slide/prose,text width=4.00cm]
      at ([xshift=0.98cm,yshift=-1.82cm]current page.north west)
      {\color{BMETerracotta}\sffamily\fontsize{6.1}{7.4}\selectfont\bfseries USEFUL FOR\par
       \vspace{0.18cm}
       \color{BMEBrown}\rmfamily\fontsize{10.0}{12.0}\selectfont\bfseries
       Local diffusion\par
       \vspace{0.20cm}
       \WarmBodyText
       Ranking, sampling, load spreading, and simple models of movement through a network.\par};

    \node[anchor=north west,inner sep=0pt,slide/prose,text width=4.00cm]
      at ([xshift=5.78cm,yshift=-1.82cm]current page.north west)
      {\color{BMESage}\sffamily\fontsize{6.1}{7.4}\selectfont\bfseries REQUIRES CARE\par
       \vspace{0.18cm}
       \color{BMEBrown}\rmfamily\fontsize{10.0}{12.0}\selectfont\bfseries
       Structure matters\par
       \vspace{0.20cm}
       \WarmBodyText
       Direction, weights, disconnected regions, and periodicity change both movement and convergence.\par};

    \node[anchor=north west,inner sep=0pt,slide/prose,text width=4.10cm]
      at ([xshift=10.62cm,yshift=-1.82cm]current page.north west)
      {\color{BMETerracotta}\sffamily\fontsize{6.1}{7.4}\selectfont\bfseries REPORT EXPLICITLY\par
       \vspace{0.18cm}
       \color{BMEBrown}\rmfamily\fontsize{10.0}{12.0}\selectfont\bfseries
       Transition rule and stopping rule\par
       \vspace{0.20cm}
       \WarmBodyText
       State whether the chain is weighted, lazy, restarted, or observed for a fixed number of steps.\par};

    \draw[BMEStone,line width=0.45pt]
      ([xshift=0.98cm,yshift=-5.23cm]current page.north west) --
      ([xshift=-0.98cm,yshift=-5.23cm]current page.north east);
    \node[anchor=north west,inner sep=0pt,slide/prose,text width=13.55cm]
      at ([xshift=1.02cm,yshift=-5.64cm]current page.north west)
      {\color{BMEBrown}\rmfamily\fontsize{11.0}{13.4}\selectfont\bfseries
       Transition rules are simple, but graph structure determines the result.\par
       \vspace{0.17cm}
       \WarmBodyText
       State the rule, its consequence, and the assumptions used to interpret the result.\par};

    \WarmFrameFooter{Random walks · modeling boundary}
  \end{tikzpicture}
\end{frame}

\end{document}
